\begin{abstract}
Laskar and Petit conjectured that a planetary system whose angular momentum deficit (AMD) is too small to allow orbits to touch is, away from short-period resonances, stable in practice. We prove a rigorous asymptotic form of the mechanism behind this conjecture, in the plane and in space, up to the sharp strong threshold, which also forbids a planet from reaching the star. Fix the mass ratios, a compact set of ordered semi-major axes and a margin $\delta>0$ below the threshold. If the planetary masses are of order $\mu$, with $\mu$ small, and the initial mean motions are Diophantine with exponent $\tau$, then for $\abs{t}\leq T_0\exp(c\mu^{-\nf{1}{(2(\tau+1))}})$ the semi-major axes move by $O(\sqrt\mu)$, the AMD stays at least $\nf{\delta}{2}$ below the threshold, and the orbital annuli stay ordered and uniformly separated. Without arithmetic conditions the same holds with Nekhoroshev's exponents. Exact conservation of angular momentum confines the AMD while the semi-major axes drift exponentially slowly; the secular dynamics may be chaotic. The constants are not explicit. The results under the Diophantine condition are formally verified in the Lean proof assistant.
\end{abstract}
